We now consider a team of $N$ robots.
Robot $i$ has state $x_i$, position $p_i=\Pi x_i\in\mathbb{R}^3$, and occupies a ball of radius $r$.
Pairwise collision avoidance between robots $i$ and $j$ is encoded by
\begin{equation}
    h_{ij}(x_i, x_j) = \|p_i - p_j \|^2 - (2r)^2.
    \label{eq:pairwise_barrier}
\end{equation}
We assume that $h_{ij}$ has relative degree one with respect to the control inputs. 
Higher-relative-degree dynamics can instead be handled using a high-order CBF construction. 
The corresponding joint CBF condition is
\begin{equation}
    L_{f_i} h_{ij} + L_{f_j}h_{ij} +
    L_{g_i} h_{ij}u_i + L_{g_j} h_{ij}u_j
    \geq -\alpha (h_{ij}).
\label{eq:centralized_CBF_cond}
\end{equation}
Because Eq.~\eqref{eq:centralized_CBF_cond} couples the two control inputs, a direct implementation jointly optimizes the stacked inputs of the robots and therefore requires coordination across the team.~\cite{wang2017safety}.

A fully unilateral alternative is to require robot $i$ to remain safe against every admissible input of robot $j$:
\begin{equation}
    \inf_{u_j\in\mathcal{U}_j} 
    \Big[ L_{f_i} h_{ij} + L_{f_j} h_{ij} + 
    L_{g_i} h_{ij} u_i + L_{g_j}h_{ij}u_j \Big] \geq -\alpha(h_{ij}).
    \label{eq:worst_case_CBF_cond}
\end{equation}
This condition provides safety independently of the neighbor's realized input, but is conservative because it neglects the neighbor's contribution and provides no mechanism for deciding which robot should yield. 
The next section addresses this limitation by distributing the shared avoidance responsibility asymmetrically.